\documentclass[11pt]{article}

\usepackage{amsmath}
\usepackage{amssymb}
\usepackage{geometry}
\usepackage{listings}
\usepackage{xcolor}

\geometry{margin=1in}

\title{Design of a Kokkos-Based CVFEM Nodal Temperature Gradient Algorithm}
\author{}
\date{}

\begin{document}

\maketitle

\section{Overview}

This document describes the design of a stand-alone
control-volume finite-element method (CVFEM) algorithm for
computing the nodal temperature gradient in a Kokkos-based
STK mesh application.

The implementation is intentionally independent of the existing
MareNalu algorithm hierarchy.  The design separates three
responsibilities:

\begin{enumerate}
\item the physics algorithm, which manages mesh parts and
orchestrates the calculation;
\item the CVFEM element implementation, which provides
element-specific operations; and
\item the Kokkos execution model, which maps mesh buckets to
Kokkos teams and elements to team threads.
\end{enumerate}

The initial implementation supports the Tet4 SCV element.
Hex8 SCV support will be added using the same interface.

\section{Algorithm Object}

A single
\texttt{ComputeNodalGradTemperatureAlg}
object represents one nodal-temperature-gradient algorithm.

The object contains references to the mesh and the fields required
by the calculation:

\begin{itemize}
\item coordinates;
\item temperature; and
\item nodal temperature gradient.
\end{itemize}

The algorithm also contains a single vector,

$$
\texttt{partVec\_},
$$

containing all mesh parts to which the algorithm applies.

Thus, multiple element parts do not require multiple algorithm
objects.  When the algorithm is registered for another part, that
part is appended to the existing \texttt{partVec\_}.

\section{Element-Type Dispatch}

The algorithm supports multiple element topologies while retaining
a single part vector.

The topology is therefore selected explicitly in
\texttt{execute()}.

For example, the locally owned element selector is first constructed
from all parts:

\begin{verbatim}
s_locally_owned_union =
locally_owned_part()
& selectUnion(partVec_);
\end{verbatim}

The selector is then restricted to a particular topology:

\begin{verbatim}
s_tet4 =
s_locally_owned_union
& get_topology_root_part(TETRAHEDRON_4);
\end{verbatim}

The Tet4 implementation is subsequently invoked through

\begin{verbatim}
execute_element<Tet4SCVElement>(s_tet4);
\end{verbatim}

The same structure will be used for Hex8:

\begin{verbatim}
execute_element<Hex8SCVElement>(s_hex8);
\end{verbatim}

Topology dispatch is deliberately visible in the main algorithm
rather than being hidden inside the templated execution routine.

\section{Nodal Output Initialization}

The nodal gradient field is an output of the algorithm.

Before the element calculations begin, the field is initialized to
zero.

The initialization selector is

\begin{verbatim}
s_all_nodes =
selectUnion(partVec_)
& selectField(gradTemperature_);
\end{verbatim}

The locally-owned restriction is deliberately not included in this
selector.  A locally owned element can contribute to a shared node,
so the nodal output must be initialized on all applicable nodes.

The output field is then made available to the device before the
device calculation.

\section{Read-Only and Read-Write Field Access}

Field access follows the STK NGP read/write convention.

The coordinates and temperature fields are dependencies of the
algorithm and are not modified:

\begin{verbatim}
auto coordinatesData =
coordinates_.data<
stk::mesh::ReadOnly,
stk::ngp::DeviceSpace>();

auto temperatureData =
temperature_.data<
stk::mesh::ReadOnly,
stk::ngp::DeviceSpace>();
\end{verbatim}

The nodal temperature-gradient field is populated by the algorithm
and is therefore accessed as read-write:

\begin{verbatim}
auto gradTemperatureData =
gradTemperature_.data<
stk::mesh::ReadWrite,
stk::ngp::DeviceSpace>();
\end{verbatim}

The distinction between read-only and read-write fields establishes
the access contract of the algorithm.

The algorithm is responsible for synchronization of the output
field that it creates and populates.  Synchronization of read-only
input fields is outside the responsibility of this algorithm.

\section{Bucket-Based Kokkos Execution}

The execution model follows the STK NGP bucket-access pattern.

The algorithm first obtains the bucket IDs associated with the
selected element topology:

\begin{verbatim}
stk::NgpVector<unsigned> bucketIds =
ngpMesh.get_bucket_ids(
stk::topology::ELEM_RANK,
selector);
\end{verbatim}

A single Kokkos \texttt{TeamPolicy} launch is then used for all
selected buckets.

The fundamental mapping is

$$
\boxed{\text{one Kokkos team} = \text{one mesh bucket}}
$$

rather than

$$
\text{one Kokkos team} = \text{one element}.
$$

Within each bucket, the elements are distributed among the team
threads using

\begin{verbatim}
Kokkos::TeamThreadRange
\end{verbatim}

This produces a nested execution structure:

$$
\text{TeamPolicy}
\longrightarrow
\text{bucket}
\longrightarrow
\text{TeamThreadRange}
\longrightarrow
\text{element}.
$$

This avoids a host-side loop over buckets followed by a separate
device kernel launch for every bucket.

\section{CVFEM Element Interface}

Element-specific CVFEM operations are contained in dedicated
element types.

The SCV and SCS element types are separate abstractions.  An SCV
element provides SCV operations, while an SCS element will provide
SCS operations.

For the current Tet4 SCV implementation, the element provides:

\begin{itemize}
\item element coordinates;
\item shape-function derivatives, \texttt{dndx};
\item SCV volumes; and
\item the mapping from SCV integration point to its associated
element node.
\end{itemize}

The element interface is represented by

\begin{verbatim}
Tet4SCVElement
\end{verbatim}

and the corresponding Hex8 implementation will be

\begin{verbatim}
Hex8SCVElement.
\end{verbatim}

\section{Element Scratch Storage}

Scratch storage is associated with the element type, but the
algorithm allocates only the portions required by the operations
that it performs.

The element therefore exposes independent scratch-size functions:

\begin{verbatim}
coordinates_scratch_size()
dndx_scratch_size()
scv_volume_scratch_size()
\end{verbatim}

There is intentionally no single aggregate
\texttt{scratch\_size()} function.

For the current nodal temperature-gradient calculation, all three
storage requirements are needed:

$$
S_{\mathrm{total}}
=
S_{\mathrm{coordinates}}
+
S_{\mathrm{dndx}}
+
S_{\mathrm{scv}}.
$$

A different algorithm that only requires coordinates and SCV volume
can allocate

$$
S_{\mathrm{coordinates}}
+
S_{\mathrm{scv}},
$$

without allocating storage for \texttt{dndx}.

This allows the element type to describe all of its available data
while allowing each algorithm to determine its actual scratch
requirements.

\section{Element Calculation}

For each element, the calculation proceeds as follows.

\begin{enumerate}
\item Obtain the connected element nodes.
\item Gather nodal coordinates from the read-only coordinates
field.
\item Gather nodal temperature from the read-only temperature
field.
\item Compute the element shape-function derivatives.
\item Compute the SCV volumes.
\item Evaluate the temperature gradient at the SCV integration
points.
\item Weight each contribution by its corresponding SCV volume.
\item Accumulate the resulting contribution into the nodal
gradient field.
\end{enumerate}

The element operations themselves are supplied by the element type:

\begin{verbatim}
Element::compute_dndx(...)
Element::compute_scv_volume(...)
Element::ip_node(...)
\end{verbatim}

Consequently, the algorithm does not contain Tet4-specific or
Hex8-specific geometric calculations.

\section{Concurrent Nodal Accumulation}

Multiple elements can contribute to the same nodal gradient.
Because elements are processed concurrently on the device, the
element contribution is accumulated using atomic operations:

\begin{verbatim}
Kokkos::atomic_add(...)
\end{verbatim}

The element first forms its local contribution and then performs
the atomic nodal accumulation.

This keeps the atomic operation outside the integration-point
calculation and limits the number of atomic updates.

\section{MPI Assembly}

The device calculation produces the local contribution to the
nodal gradient.

After the device calculation is complete, the output field is
synchronized to the host and STK's parallel reduction is used to
combine contributions to shared nodes:

\begin{verbatim}
stk::mesh::parallel_sum(
bulk_,
{&gradTemperature_});
\end{verbatim}

The overall operation is therefore

$$
\boxed{
\text{initialize}
\rightarrow
\text{device element calculation}
\rightarrow
\text{host synchronization}
\rightarrow
\text{MPI sum}
}
$$

\section{Overall Execution Structure}

The complete algorithm has the following conceptual structure:

\begin{verbatim}
execute()
|
+-- construct locally-owned element selector
|
+-- construct nodal output selector
|
+-- zero gradTemperature
|
+-- synchronize gradTemperature to device
|
+-- select Tet4 elements
|     |
|     +-- execute_element<Tet4SCVElement>()
|            |
|            +-- obtain bucket IDs
|            +-- one Kokkos team per bucket
|            +-- TeamThreadRange over elements
|            +-- element CVFEM operations
|            +-- atomic nodal accumulation
|
+-- select Hex8 elements
|     |
|     +-- execute_element<Hex8SCVElement>()
|            |
|            +-- same execution structure
|
+-- synchronize gradTemperature to host
|
+-- MPI parallel_sum
\end{verbatim}

\section{Design Principles}

The implementation follows these principles:

\begin{enumerate}
\item One algorithm object can operate on multiple mesh parts.
\item A single \texttt{partVec\_} is used for all supported
element types.
\item Topology dispatch is explicit in \texttt{execute()}.
\item One Kokkos team represents one mesh bucket.
\item Elements within a bucket are processed with
\texttt{TeamThreadRange}.
\item Element operations are separated from algorithm orchestration.
\item SCV and SCS element types remain distinct.
\item Scratch requirements are exposed independently for each
element operation.
\item Read-only and read-write field access is explicitly declared.
\item The algorithm owns synchronization of the output field that
it creates and populates.
\item Read-only input-field synchronization is not owned by this
algorithm.
\item MPI assembly is performed after the device calculation.
\end{enumerate}

\section{Initial Implementation Scope}

The initial implementation is intentionally limited to Tet4 SCV
operations:

$$
\text{Tet4 SCV}
\quad\rightarrow\quad
\left\{
\texttt{dndx},
\texttt{scv\_volume}
\right\}.
$$

Hex8 SCV support will subsequently implement the same interface.

SCS operations are outside the scope of the current nodal
temperature-gradient calculation and will be introduced through a
separate SCS element type.

\end{document}
